\section{The Communion of Saints, Our Lady Queen of Heaven, and the Angels}

Sacred Scripture never depicts heaven as a solitary desert where isolated souls hover before an impersonal Absolute. Its constant figures are a household, a wedding banquet, a liturgy, and a city. The Epistle to the Hebrews lifts believers' eyes to the assembly they have already entered by faith and baptism:

\begin{quote}
But you are come to mount Sion and to the city of the living God, the heavenly Jerusalem, and to the company of many thousands of angels, and to the church of the firstborn who are written in the heavens, and to God the judge of all, and to the spirits of the just made perfect, and to Jesus the mediator of the new testament, and to the sprinkling of blood which speaketh better than that of Abel. (Heb 12:22--24)
\end{quote}

St.\ John beholds that same living company in Apocalypse 7: ``a great multitude, which no man could number, of all nations and tribes and peoples and tongues, standing before the throne and in sight of the Lamb, clothed with white robes, and palms in their hands'' (Apoc 7:9), joined by all the angels, ancients, and living creatures in one chorus: ``Amen. Benediction and glory and wisdom and thanksgiving, honour and power and strength, to our God, for ever and ever. Amen'' (Apoc 7:11--12). Every redeemed soul enters the \textbf{Communion of Saints} perfected---the Church Triumphant united with the angelic host around Christ the Head and Mary His Mother.

\subsection*{The Nine Choirs of Angels}

Before any human foot entered heaven, the celestial city was already ablaze with pure spirits who stood steadfast when Lucifer fell. Preaching in Rome on the parables of the lost sheep and the ten drachmas (Luke 15:1--10), St.\ Gregory the Great weaves together the whole biblical testimony on the angels, teaching that the woman possessed ten drachmas because there are nine orders of angels, and man was created as the tenth to complete the number of the elect (\textit{Decem vero drachmas habuit mulier, quia novem sunt ordines angelorum. Sed ut compleretur electorum numerus, homo decimus est creatus}).\endnote{St.\ Gregory the Great, \textit{Homilies on the Gospels} (\textit{Homiliae in Evangelia}) 34.6--10 (PL 76).}

Gathering the scriptural names from Isaias, Ezechiel, Ephesians, Colossians, and Jude, Gregory enumerates the nine choirs---Angels, Archangels, Virtues, Powers, Principalities, Dominations, Thrones, Cherubim, and Seraphim (\textit{angelos, archangelos, virtutes, potestates, principatus, dominationes, thronos, cherubim, atque seraphim}).\endnote{St.\ Gregory the Great, \textit{Homiliae in Evangelia} 34.7; cf.\ Isa 6:2; Ezech 1 and 10; Eph 1:21; Col 1:16; Jude 1:9.} \textit{Angel}, he notes, names an office, not a nature: those blessed spirits are always spirits, yet called angels when sent as messengers---Michael (``Who is like God?''), Gabriel (``Strength of God''), Raphael (``Medicine of God'') (\textit{Hom.}\ 34.8--9)---ascending through Virtues, Powers, Principalities, Dominations, and Thrones to the Cherubim, filled with knowledge, and the Seraphim, burning with love.

Following Pseudo-Dionysius and St.\ Gregory in \textit{ST} I, q.\ 108, aa.\ 1 and 5, St.\ Thomas Aquinas sets these nine choirs in three hierarchies of three orders each, noting that Dionysius explains the names from each choir's interior perfection and Gregory from its outward ministry.\endnote{St.\ Thomas Aquinas, \textit{Summa Theologiae} I, q.\ 108, aa.\ 1, 5.} Yet angels and redeemed men form one single commonwealth under God. Quoting St.\ Augustine (\textit{City of God} XII.1), Aquinas writes: ``There are two cities, that is, two societies, one of the good angels and men, the other of the wicked'' (\textit{ST} I, q.\ 108, a.\ 1). When the Good Shepherd carries the lost sheep home, Gregory preaches, He calls His ``friends and neighbours''---the faithful angels who continually behold His face---to rejoice (\textit{Hom.}\ 34.3), welcoming redeemed men as their fellow-citizens (\textit{City of God} XXII.29).

\subsection*{Our Lady, Assumed and Crowned Queen of Heaven}

At the summit of all creation---below the uncreated Trinity alone, yet exalted above every choir of Cherubim and Seraphim---stands the Blessed Virgin Mary, Mother of God. In Apocalypse 12:1, St.\ John sees a sign applied by Catholic liturgy and tradition at once to the Church and pre-eminently to Mary: ``And a great sign appeared in heaven: A woman clothed with the sun, and the moon under her feet, and on her head a crown of twelve stars.''

Two solemn acts of Pope Pius XII gathered the Church's immemorial faith into definitive and universal teaching. On All Saints' Day, 1 November 1950, in the Apostolic Constitution \textit{Munificentissimus Deus}, Pius XII exercised his supreme apostolic authority to pronounce, declare, and define \textit{ex cathedra} as a divinely revealed dogma that the Immaculate Mother of God, the ever-Virgin Mary, having completed the course of her earthly life, was assumed body and soul into heavenly glory.\endnote{Pope Pius XII, Apostolic Constitution \textit{Munificentissimus Deus} (1 November 1950), \S 44 (DS 3903).} Preserved from original sin and united wholly to the saving Passion and victory of her Son, Mary does not wait for the Last Day to resume her body; she already stands in heaven in the full integrity of body and soul, the firstfruits and pledge of the Church's resurrection.

Four years later, on 11 October 1954, Pius XII issued the encyclical \textit{Ad Caeli Reginam}, proclaiming Mary's royal dignity as Queen of Heaven and instituting her universal feast.\endnote{Pope Pius XII, Encyclical Letter \textit{Ad Caeli Reginam} (11 October 1954).} Grounding her queenship on her Divine Motherhood---by which she is uniquely joined to Christ, true God and eternal King---and on her singular association as the New Eve in His redemptive Passion, the Pope shows that her throne at her Son's right hand is a queenship of maternal intercession and mercy. St.\ Alphonsus Liguori captures the tenderness of that welcome in Consideration XXIX of \textit{Preparation for Death}: entering Paradise, ``the soul shall see all the graces she obtained through the intercession of Mary, who shall embrace her with love and tenderness. This queen of Heaven shall then conduct the soul to Jesus, who will receive her as his spouse.''\endnote{St.\ Alphonsus Maria de Liguori, \textit{Preparation for Death}, Consideration XXIX, First Point (Dublin: James Duffy, 1843), pp.\ 351--352.}

\subsection*{The Fellowship of the Saints}

Does a soul possessing the infinite God need human friends to be happy? In \textit{ST} I-II, q.\ 4, a.\ 8, St.\ Thomas Aquinas keeps both divine sufficiency and human charity in exact balance. Because man has the entire fullness of his perfection in God, ``the fellowship of friends is not essential to Happiness''; ``if there were but one soul enjoying God, it would be happy, though having no neighbor to love.'' Yet fellowship belongs genuinely to the \textit{well-being} (\textit{bene esse}) of heaven:

\begin{quote}
But the fellowship of friends conduces to the well-being of Happiness. Hence Augustine says (\textit{Gen.\ ad lit.}\ viii, 25) that ``the spiritual creatures receive no other interior aid to happiness than the eternity, truth, and charity of the Creator. But if they can be said to be helped from without, perhaps it is only by this that they see one another and rejoice in God, at their fellowship.'' \dots\ Supposing one neighbor to be there, love of him results from perfect love of God. Consequently, friendship is, as it were, concomitant with perfect Happiness.\endnote{St.\ Thomas Aquinas, \textit{Summa Theologiae} I-II, q.\ 4, a.\ 8, corpus and ad 3.}
\end{quote}

Every clean human bond formed in Christ is not erased at heaven's gate, but purified and restored in God. The \textit{Roman Catechism} counts this mutual recognition among the principal accessory joys of eternal life: ``that glory which consists in the clear and comprehensive knowledge, which each of the blessed shall have of the singular and exalted dignity of his companions in glory.''\endnote{\textit{Catechism of the Council of Trent}, Part I, Creed Article XII, trans.\ Donovan, pp.\ 96--97.} Nor do the blessed forget the Church Militant on earth. Citing St.\ Gregory the Great, the \textit{Supplement} (q.\ 94, a.\ 1, ad 1) affirms that ``the saints in heaven know distinctly all that happens both to wayfarers and to the damned,'' because ``since they possess the glory of God within them, we cannot believe that external things are unknown to them.''\endnote{\textit{Supplement}, q.\ 94, a.\ 1, ad 1 (English Dominicans, 1922).} Beholding us in the Word of God, they lift our prayers before the throne like golden vials of incense (Apoc 5:8; CCC 1029).

St.\ Francis de Sales gathers the warmth of this celestial company into his Eighth Meditation in \textit{Introduction to a Devout Life}:

\begin{quote}
Consider the glory, the beauty, and the multitude of the inhabitants of this happy country; millions of millions of angels, of cherubim and seraphim; choirs of apostles, prophets, martyrs, confessors, virgins, and holy women; the multitude is innumerable. O how glorious is this company! the least of them is more beautiful to behold than the whole world; what a sight then will it be to behold them all! But, O my God! how happy are they! they sing incessantly harmonious songs of eternal love! they always enjoy a state of felicity; they mutually give each other unspeakable contentment, and live in the consolation of a happy, indissoluble society.\endnote{St.\ Francis de Sales, \textit{Introduction to a Devout Life} I.16 (Eighth Meditation: On Heaven), Roman Catholic Pustet ed., pp.\ 40--41.}
\end{quote}

``They are like happy birds,'' Francis adds in the next paragraph, ``flying and singing perpetually in the air of his divinity, which encompasses them on all sides with incredible pleasure. There every one does his utmost, and sings without envy the praises of his Creator.''
